\section{Statement of Patient Commitment}

The parent clinical protocol opens with a Statement of Compliance addressed to
the Food and Drug Administration, the Institutional Review Board, and the
sponsor \cite{onpremwhippl}. That statement is necessary and it is not
addressed to the person whose abdomen will be opened. This paper opens instead
with the seven things the same protocol commits to the participant, each one
tied to the clause that makes it enforceable rather than aspirational.

The premise is taken from the author's earlier analysis of proposed United
States legislation for Physical AI oncology trials: the cancer patient, not the
doctor, the nurse, the trial sponsor, the Institutional Review Board, or the
regulator, is the priority participant \cite{paibillprior, hr9510billv5}.

\draftinstr{Expand this opening to roughly 1,300 characters. Source the framing
from \texttt{inputs/patient-priority-physical-ai.zip} and cite
\texttt{hr9510billv5} for the legislative premise. State explicitly that this is
an advocacy document written from the protocol's own record, not a substitute
for the consent form.}

\subsection{The seven commitments}

\draftinstr{Write one short paragraph per commitment, roughly 420 characters
each, in the order below. Each paragraph names the commitment, states the
protocol clause that carries it, and states what a participant can check to
verify it. Source every clause from
\texttt{inputs/phase-1-trial-protocol.zip}, sections \texttt{sec-00-compliance},
\texttt{sec-04-design}, and \texttt{sec-06-intervention}.}

\vspace{0.30em}

\noindent\begin{tabularx}{\textwidth}{L{0.7cm} L{4.6cm} Y}
\toprule
\textbf{\#} & \textbf{Commitment} & \textbf{The clause that makes it enforceable} \\
\midrule
1 & Nothing starts without you & Consent under 21 CFR \S 50.20 precedes every study procedure \\
2 & The robot is proven before you & Phase 0 gate: at least 1000 simulations across at least two independent frameworks, Unified Safety Level at least 7.0 \\
3 & You are never the untested dose & 3+3 escalation with sentinel staggering; a dose level opens only after the level below is cleared \\
4 & A human approves every motion & Class II collaborative device under continuous human oversight; no model-to-actuator path exists \\
5 & Stopping is faster than harm & Emergency stop at most 3 ms across arms and at most 500 ms system-wide \\
6 & Your record cannot be edited & Hash-chained audit trail under 21 CFR part 11, bound to a deterministic seed and a repository commit \\
7 & Leaving costs you nothing & Withdrawal at any visit without reason; standard care continues unchanged \\
\bottomrule
\end{tabularx}

\vspace{0.30em}

\draftinstr{Table above is final in structure. Verify the column widths sum to
\texttt{textwidth}: 0.7cm plus 4.6cm plus the Y remainder, with
\texttt{tabcolsep} 5pt on six edges.}

\subsection{What this paper is, and what it is not}

\draftinstr{Roughly 900 characters. It is an independent advocacy paper and a
practical adoption guide built from the author's repository sources. It is not
the consent form, not medical advice, not regulatory advice, and not endorsed by
any sponsor, CRO, site, IRB, regulator, or medical society. Cite
\texttt{NCIConsent24} for what a consent form is required to contain, and state
that this paper is designed to be read before that form, not instead of it.}

\subsection{How to read the thirty figures}

\draftinstr{Roughly 800 characters. Explain the five diagram types and the
question each answers, drawing the wording from the five stage READMEs at
\texttt{mermaid/README.md}, \texttt{plantuml/README.md}, \texttt{d2/README.md},
\texttt{diagrams-python/README.md}, and \texttt{graphviz/README.md}. State that
the monospace tag at the top left of every frame names the native construct the
figure reproduces, so a reader can carry any figure back to its platform.}

\begin{pafig}[mermaid-type: flowchart TD]
\node[mmgoal] (a) at (0,0) {Figure 1 slot\\seven commitments\\flowchart TD};
\node[mmin] (b) at (0,-2.2) {Source:\\mermaid/fig-01-seven-commitments.md};
\draw[mmedge] (a) -- (b);
\end{pafig}
\vspace{-0.7cm}
\figcaption{Figure 1. The seven commitments made to the participant, from the first
conversation through consent, the Phase 0 gate, and dose assignment, to the day-30\\
safety window and the day-90 pathology result that is reported back to them directly.}

\draftinstr{Replace the Figure 1 slot with the full mermaid-type flowchart
specified in \texttt{mermaid/fig-01-seven-commitments.md}, following the TikZ
rendering notes at the foot of that file: seven \texttt{mmin} commitment nodes
in one column at 2.4 cm pitch, three \texttt{mmdec} gates at $x=5.6$, one
\texttt{mmdark} halt node reached by three \texttt{mmedged} edges, and the
day-30 and day-90 nodes below.}

\subsection{Accountability, stated once, at the front}

\draftinstr{Roughly 700 characters. The single most common accountability
question in the surveyed literature is "who is to blame?"
\cite{Jauniaux2025, BrarRAS2024X}. Answer it here in summary and forward to
\S11: there is exactly one accountable party per failure class, named in the
responsibility matrix of Figure 29, and the participant is told which one to
call.}

\begin{pafig}[graphviz-type: digraph, rooted DAG]
\node[gvkey] (a) at (0,0) {Figure 2 slot\\accountability DAG};
\node[gvbox] (b) at (0,-2.0) {Source:\\graphviz/fig-02-accountability-dag.dot};
\draw[gvedge] (a) -- (b);
\end{pafig}
\vspace{-0.7cm}
\figcaption{Figure 2. The accountability chain drawn as a rooted directed acyclic graph
with the participant at the root, showing the people in the room, the bodies accountable\\
for the study, the independent reviewers, and the three obligations that point back.}

\draftinstr{Replace the Figure 2 slot with the full graphviz-type DAG specified
in \texttt{graphviz/fig-02-accountability-dag.dot}. Preserve the stated
invariant: the graph is acyclic, so no body in the chart is its own reviewer.
Draw exactly three back-reference edges to the root, no more, because three is
what the protocol actually owes the participant.}
